%% =====================================================================
%%  RESUME (2 pages) -- Higher education: teaching-focused faculty,
%%  learning / quantitative center leadership, academic support
%%  Wöden Kusner
%% =====================================================================
%%  Build: pdflatex Kusner_Resume_Higher-Ed-Teaching-Support.tex (twice)
%%
%%  One column, plain headings, real text, dates on the role line, so
%%  applicant-tracking systems (Workday, Interfolio, PeopleAdmin) parse it.
%%  Tailor the SUMMARY to each posting; keep the QRC section first.
%% =====================================================================
\documentclass[11pt]{article}
\usepackage[letterpaper,margin=0.7in]{geometry}
\usepackage[expansion=false]{microtype}
\input{glyphtounicode}\pdfgentounicode=1
\usepackage{enumitem}
\usepackage{xcolor}
\usepackage[hidelinks]{hyperref}
\hypersetup{pdftitle={Woden Kusner - Resume}, pdfauthor={W\"oden Kusner (Woden Kusner)},
  pdfkeywords={Woden Kusner, mathematics, quantitative reasoning, tutoring center, learning assistants,
  peer tutoring, academic support, student success, ALEKS placement, teaching, mentoring}}
\pagestyle{empty}
\setlength{\parindent}{0pt}
\definecolor{rule}{gray}{0.55}
\hyphenpenalty=10000 \exhyphenpenalty=10000 \emergencystretch=3em  % no split words: cleaner ATS parsing

% --- helpers -----------------------------------------------------------
\newcommand{\sect}[1]{\par\vspace{7pt}{\large\bfseries #1}\par\vspace{1pt}%
  {\color{rule}\hrule height .5pt}\vspace{4pt}}
\newcommand{\role}[3]{\par\vspace{3pt}\textbf{#1}\ \hfill\mbox{#3}\par   % title, org, dates
  {\itshape #2}\par}
\setlist[itemize]{leftmargin=1.1em, itemsep=1pt, topsep=1pt, parsep=0pt, label=$\bullet$}

\begin{document}

%% ---------------- HEADER ----------------
{\LARGE\bfseries W\"oden Kusner}\par\vspace{2pt}
Mathematics educator $|$ quantitative support and peer-learning programs\par\vspace{2pt}
\href{mailto:wkusner@gmail.com}{wkusner@gmail.com} $|$
\href{https://wkusner.github.io}{wkusner.github.io} $|$
Appleton, WI

%% ---------------- SUMMARY ----------------
\sect{Summary}
Mathematician and teacher (Ph.D., University of Pittsburgh) with teaching experience at five universities since
2009, now directing the Quantitative Resource Center at Lawrence University. Builds peer-learning programs,
supports students who find mathematics difficult, and works across departments so that quantitative courses
and student support fit together. Holds certificates in college teaching and in diversity and inclusion.

%% ---------------- STRENGTHS ----------------
\sect{Areas of Strength}
\begin{itemize}
  \item Peer tutoring and learning-assistant programs: design, hiring, training, mentoring, and scheduling.
  \item Individualized support for students, including case-managed students; open office hours for
        students, faculty, and staff.
  \item Course design for introductory mathematics; ALEKS placement and adaptive-learning courses.
  \item Teaching probability, statistics, and calculus from preparatory through advanced undergraduate levels.
  \item Mentoring undergraduate research, directed studies, and graduate-school preparation.
\end{itemize}

%% ---------------- EXPERIENCE ----------------
\sect{Experience}
\role{Director, Quantitative Resource Center; Adjunct Assistant Professor of Mathematics}{Lawrence University}{2025--present}
\begin{itemize}
  \item Direct quantitative peer tutoring for the college: recruit, train, and mentor peer tutors and coordinate
        support with departments.
  \item Created the Quantitative Learning Assistant program (2025): wrote the position description; hire and
        schedule specialized tutors, each paired with a course and its instructor.
  \item Provide individualized support for case-managed students; hold open office hours for students,
        faculty, and staff; advise on ALEKS math placement.
  \item Presented on the Quantitative Resource Center and perspectives on mathematics at a Center for Teaching
        Excellence Lunch and Learn (2025); part of Lawrence's team at a Carleton College workshop on
        institutional approaches to quantitative skill support (2025).
  \item Teach introductory mathematics (Foundations in Math, Preparation for Calculus, Calculus) and supervise
        directed studies.
\end{itemize}
\role{Instructor of Mathematics}{Lawrence University}{2023--2025}
\begin{itemize}
  \item Coordinated peer tutoring and math support for introductory mathematics.
  \item Designed the ALEKS-based course structure and exercise pack for Preparation for Calculus, later used by
        colleagues teaching the course; supervised undergraduate teaching assistants.
  \item Mentored a summer undergraduate research fellow (discrepancy theory) and ran a graduate-school
        preparation program (GRE mathematics practice and applications) for three students.
\end{itemize}
\role{Assistant Professor (non-tenure-track)}{University of Georgia}{2020--2022}
\begin{itemize}
  \item Taught Calculus I and II in a coordinated multi-section program; served on the calculus exam committee.
  \item UGA Teaching Academy Fellow; co-organized the Math Learning and Teaching Seminar; member of the
        instructional corps recognized by the 2022 Regents' Teaching Excellence Award.
\end{itemize}
\role{Assistant Professor (non-tenure-track) and Postdoctoral Associate}{Vanderbilt University}{2017--2020}
\begin{itemize}
  \item Taught probability and statistical inference, mathematical statistics, and multivariable calculus;
        supervised and mentored teams of undergraduate teaching assistants and graders.
  \item Mentored undergraduate research (including a Highest Honors thesis) and served on honors and
        dissertation committees.
\end{itemize}
\role{FWF Postdoctoral Researcher and Lecturer}{Graz University of Technology, Austria}{2014--2017}
\begin{itemize}
  \item Taught a course on packings, lattices, and configurations; assisted in higher analysis.
\end{itemize}
\role{Teaching Fellow and Teaching Assistant}{University of Pittsburgh}{2009--2011}
\begin{itemize}
  \item Led recitations for calculus, business calculus, foundations, algebra, and topology courses; taught
        Business Calculus as instructor of record.
\end{itemize}

%% ---------------- EDUCATION ----------------
\sect{Education and Certificates}
\textbf{Ph.D., Mathematics}, University of Pittsburgh \hfill 2014\par
\textbf{B.S., Mathematics}, Haverford College \hfill 2007\par\vspace{3pt}
Certificate in Diversity and Inclusion, University of Georgia \hfill 2021\par
STEM Teaching Specialization, Vanderbilt University \hfill 2020\par
Certificate in College Teaching, Vanderbilt University \hfill 2019\par
Community Engaged Learning faculty cohort, Lawrence University \hfill 2026--2027\par
REFLECT peer observation of teaching, Lawrence University \hfill 2024

%% ---------------- SERVICE / RESEARCH ----------------
\sect{Mentoring, Service, and Scholarship}
\begin{itemize}
  \item Letters of recommendation for more than 20 students (2018--2025): graduate school, research programs,
        medical school, transfer and study abroad, employment.
  \item LS-PAC MODELS mentor (2021--2023); leadership committee, Haverford Multicultural Alumni Action
        Group (2020--2024).
  \item Community: Historic Central Neighborhood Association (Appleton), Safe Routes to School subcommittee,
        Fox Valley Data Exchange working group.
  \item Research in discrete geometry: 7 peer-reviewed papers, 37 talks; referee for \emph{Annals of
        Mathematics} and other journals; NSF Graduate Research Fellowship reviewer (2022).
\end{itemize}

\end{document}
